\section{Bosonic states, energy, and entropy}\label{sec:preliminaries}
We work on the \(n\)-mode Fock space
\(\mathcal H_n=\ell^2(\N_0^n)\), with number basis
\(\{|\mathbf k\rangle:\mathbf k\in\N_0^n\}\). For the annihilation operators \(d_j\),
\[
 d_j|\mathbf k\rangle=\sqrt{k_j}\,|\mathbf k-\mathbf e_j\rangle,
 \qquad N=\sum_{j=1}^n d_j^\dagger d_j,
 \qquad W=\id+N.
\]
A state is a positive trace-class operator of trace one. Expectations of positive unbounded operators are understood as increasing limits of bounded spectral truncations. Form inequalities have their usual quadratic-form meaning on the indicated domains.

For \(r>0\), put
\begin{equation}\label{eq:thermal-definition}
 h(r)=\log(1+1/r),\qquad
 \tau_r=(r+1)^{-n}e^{-h(r)N}.
\end{equation}
The product geometric law gives
\begin{equation}\label{eq:thermal-moments}
 \Tr\tau_rN=nr,\qquad S(\tau_r)=ng(r).
\end{equation}
The state \(\tau_r\) is faithful and has all number moments finite. At \(r=0\), its limiting state is the multimode vacuum.

\subsection{Entropy estimates and compactness}
We first record the entropy facts needed in the infinite-dimensional argument.
The relative entropy introduced by Umegaki \cite{Umegaki1962} gives the
Gibbs variational principle. For oscillator energy bounds, the compactness
and entropy-continuity statements below are special cases of the general
energy-constrained theory \cite{Shirokov2006,Winter2016}; we include the
elementary cutoff proof used here.

\begin{lemma}[Relative entropy and the Gibbs variational inequality]\label{lem:gibbs-variational}
Let \(\sigma\) be a faithful state, and suppose both \(S(\rho)\) and \(\Tr\rho(-\log\sigma)\) are finite. Then
\[
 D(\rho\Vert\sigma):=-S(\rho)+\Tr\rho(-\log\sigma)\ge0,
\]
with equality if and only if \(\rho=\sigma\).
Consequently, if \(K\) is bounded below and \(0<Z=\Tr e^{-K}<\infty\), its Gibbs state \(Z^{-1}e^{-K}\) uniquely minimizes \(\Tr\rho K-S(\rho)\) among states for which these expressions are finite, provided that Gibbs state belongs to the admitted class.
\end{lemma}
\begin{proof}
Choose eigenvectors \(\psi_j\) of \(\rho\) with positive eigenvalues \(p_j\), and set \(q_j=\langle\psi_j,\sigma\psi_j\rangle\). Scalar Jensen applied to the spectral measure of \(\sigma\) gives
\[
 -\langle\psi_j,\log\sigma\,\psi_j\rangle\ge-\log q_j.
\]
The elementary inequality \(p\log(p/q)\ge p-q\) and \(\sum_jq_j\le1\) now imply nonnegativity. Equality requires \(p_j=q_j\), no mass of \(\sigma\) outside the support of \(\rho\), and equality in each Jensen inequality. The last condition makes each \(\psi_j\) an eigenvector of \(\sigma\) with eigenvalue \(q_j\), so \(\rho=\sigma\). The Gibbs assertion follows by substituting \(-\log\sigma=K+\log Z\).
\end{proof}

\begin{lemma}[Energy bound and compactness]\label{lem:energy-compactness}
For each finite \(E\ge0\), the set of states with \(\Tr\rho N\le E\) is compact in trace norm. On this set entropy is continuous and obeys
\begin{equation}\label{eq:entropy-energy}
 0\le S(\rho)\le n g(E/n).
\end{equation}
Every finite-energy state is a trace-norm limit of normalized finite-number truncations with uniformly bounded energies and convergent entropies.
\end{lemma}
\begin{proof}
Diagonal entropy in any orthonormal basis bounds von Neumann entropy above: apply concavity of \(-t\log t\) to each diagonal entry expressed in a spectral decomposition, then sum the nonnegative terms. Compare the number-basis probabilities with the geometric probabilities of \(\tau_{E/n}\). Summing
\(-p\log p\le-p\log q+q-p\) gives, when \(E>0\),
\[
 S(\rho)\le n\log(1+E/n)+h(E/n)E=ng(E/n).
\]
For \(E=0\), the state is the vacuum and the assertion is immediate.

Let \(P_K=\mathbf1_{\{N\le K\}}\), \(Q_K=\id-P_K\), and \(q=\Tr\rho Q_K\). The rank of \(P_K\) is finite, and \(q\le E/(K+1)\). Factoring off \(\rho^{1/2}\) and applying the Hilbert--Schmidt Cauchy--Schwarz inequality bounds each off-diagonal block by \(\sqrt q\) in trace norm. Hence \(\rho\) is uniformly approximated by its finite-dimensional compression; normalizing the latter changes its trace norm by at most \(q\). This proves precompactness. The energy bound is closed by lower semicontinuity of positive spectral expectations, so the set is compact.

For completeness, write \(H_2(q)=-q\log q-(1-q)\log(1-q)\). Pinching into the two number blocks changes the entropy by a quantity between zero and \(H_2(q)\). To see the upper bound directly, split each vector in a pure eigenensemble of \(\rho\) into its two projected vectors. Concavity of \(H_2\) bounds the increase of the ensemble-weight entropy by \(H_2(q)\); the entropy of a pure-state mixture is at most its weight entropy. The latter fact follows by applying diagonal entropy to the Gram operator of the weighted vectors, which has the mixture's nonzero eigenvalues. The lower bound follows from the same concavity argument for pinching.

The pinched entropy equals
\[
 H_2(q)+(1-q)S\!\left(\frac{P_K\rho P_K}{1-q}\right)
       +qS\!\left(\frac{Q_K\rho Q_K}{q}\right),
\]
where zero-mass terms are omitted. The last term is bounded by
\[
 qn\,g\!\left(\frac{E}{qn}\right),
\]
which tends to zero uniformly as \(q\downarrow0\). The low-number entropy term is continuous in the finite-dimensional compression. These bounds uniformly approximate \(S(\rho)\) by continuous finite-dimensional expressions and prove its continuity under the common energy bound.

Finally, normalized number compressions converge in trace norm. Their energies are at most \(E/(1-q)\), and hence are uniformly bounded for all sufficiently large \(K\). The established continuity proves entropy convergence.
\end{proof}

In particular, entropy grows only logarithmically with energy for fixed \(n\), whereas a positive linear energy penalty is coercive. For \(r>0\), we shall repeatedly use the identity
\begin{equation}\label{eq:thermal-relative-entropy}
 D(\rho\Vert\tau_r)=-S(\rho)+n\log(r+1)+h(r)\Tr\rho N.
\end{equation}

\subsection{Passive mixing}
For any finite collection of ports, a real orthogonal rotation of their mode operators, performed identically for every mode index, is implemented by a unitary on their joint Fock space. One way to construct it is to rotate the creation operators acting on the joint vacuum; the commutation relations show that the resulting number vectors again form an orthonormal basis. This unitary preserves combined total number. A retained subsystem therefore has energy no greater than the combined input energy. On each finite total-number sector, the unitary depends continuously on the rotation; consequently the induced channels are trace-norm continuous on every fixed state as the rotation varies.

\begin{proposition}[Thermal mixing]\label{prop:thermal-mixing}
Suppose independent ports carry the \(n\)-mode states \(\tau_{r_s}\), and the retained modes are \(\sum_s\sqrt{\lambda_s}\,d_{s,j}\), where \(\lambda_s\ge0\) and \(\sum_s\lambda_s=1\). Their joint retained state is
\[
 \tau_{\sum_s\lambda_s r_s}.
\]
\end{proposition}
\begin{proof}
The coherent-state representation of a thermal state is
\[
 \tau_r=\int_{\C^n}|z\rangle\langle z|\,
       \frac{e^{-\lVert z\rVert^2/r}}{(\pi r)^n}\,d^{2n}z.
\]
Indeed, inserting the coefficients
\(\langle\mathbf k|z\rangle=e^{-\lVert z\rVert^2/2}z^{\mathbf k}/\sqrt{\mathbf k!}\)
and integrating gives the diagonal geometric probabilities in \eqref{eq:thermal-definition}. Products of coherent states transform by the same linear rotation of their amplitudes, as follows from their annihilation eigenvector equations or creation-operator expansion. Independent centered circular complex Gaussians with covariances \(r_s\id\) consequently produce retained covariance \((\sum_s\lambda_s r_s)\id\). The displayed representation proves the claim. Zero reference means follow by continuity.
\end{proof}

This thermal representation and its relation to entropy photon number are also the starting point of the original formulation \cite[Section~II.B]{GES2007}. In the argument below it is only the auxiliary reference states that have this special form.
